\chapter{Client-Server Interaction Flow}

This section details the ``Handshake'' lifecycle when a user searches for a route.

\section{Phase 1: Coordinate Resolution (Client-Side)}

The coordinate resolution phase begins when a user enters a location query and concludes when the system presents selectable location results.

\begin{tcolorbox}[colback=gray!10,colframe=black!50,arc=2mm,title=Coordinate Resolution Flow]
\begin{verbatim}
User Input: "Ben Thanh Market"
    ↓
RoutesPage.js calls api.searchLocation("Ben Thanh Market")
    ↓
api.js constructs POST request to /search
    ↓
Backend returns: [
  {
    "address": {"freeformAddress": "Ben Thanh Market, District 1"},
    "position": {"lat": 10.7721, "lon": 106.6983}
  }
]
    ↓
React State updated: selectedLocation = {
  name: "Ben Thanh Market",
  lat: 10.7721,
  lon: 106.6983
}
    ↓
Dropdown displays search results
\end{verbatim}
\end{tcolorbox}

\section{Phase 2: Route Calculation (The Critical Handshake)}

The route calculation phase represents the core interaction between client and server, where the system processes geographic coordinates and returns navigable route information.

\begin{tcolorbox}[colback=gray!10,colframe=black!50,arc=2mm,title=Route Calculation Flow]
\begin{verbatim}
User clicks "Find Route"
    ↓
RoutesPage.js calls api.calculateRoute(start, end, travelMode)
    ↓
api.js constructs POST request to /route
    ↓
Request: {
  "start": {"lat": 10.7721, "lon": 106.6983},
  "end": {"lat": 10.7621, "lon": 106.6869},
  "travelMode": "car"
}
    ↓
Backend processes:
1. Validates coordinates are within HCM City
2. Queries routing engine (TomTom API)
3. Calculates distance and time
4. Returns polyline coordinates
    ↓
Response: {
  "coords": [...],
  "distance_km": 12.5,
  "duration_min": 25.3
}
    ↓
React State updated: coords = [...]
    ↓
Maps.js receives coords
    ↓
React-Leaflet renders <Polyline /> with coordinates
    ↓
Blue line appears on map
\end{verbatim}
\end{tcolorbox}

\section{Phase 3: Rendering (Desired State)}

The rendering phase transforms backend data into visual components that provide users with comprehensive route information and interactive map elements.

\begin{tcolorbox}[colback=gray!10,colframe=black!50,arc=2mm,title=Rendering Flow]
\begin{verbatim}
Backend returns JSON object
    ↓
RoutesPage receives data
    ↓
Passes to Maps.js component
    ↓
Maps.js renders:
1. Start marker (green)
2. End marker (red)
3. Route polyline (blue)
    ↓
Simultaneously, SearchBoxRoutes displays:
- Distance: 12.5 km
- Duration: 25 minutes
- Transport mode: Car
    ↓
User sees complete route visualization
\end{verbatim}
\end{tcolorbox}

\subsection{Key Interaction Points}

The client-server interaction involves several critical handshake points:

\begin{itemize}
    \item \textbf{Location Search:} User input is transformed into structured geographic data through the \texttt{/search} endpoint.
    \item \textbf{Route Request:} Coordinate pairs and travel mode preferences are transmitted to the backend for processing.
    \item \textbf{Data Validation:} The backend ensures coordinates fall within the HCM City geographic bounds before querying external routing services.
    \item \textbf{Polyline Rendering:} The frontend receives coordinate arrays and delegates rendering to the React-Leaflet mapping library.
    \item \textbf{State Synchronization:} React state management ensures UI components reflect the current route calculation results.
\end{itemize}
